\section{Analog side}

\begin{frame}{Analog flow: topology $\rightarrow$ sizing $\rightarrow$ SPICE}
  \begin{center}
  \begin{tikzpicture}[node distance=0.8cm]
    \node[agent]  (nl)  {natural-language\\spec};
    \node[agent, right=of nl]  (rec) {analog-topology\\recommender};
    \node[agent, right=of rec] (adv) {analog-sizing\\advisor};
    \node[agent2, below=0.9cm of adv] (ar)  {AutoresearchRunner\\\scriptsize greedy gm/ID loop};
    \node[agent3, left=of ar] (ng) {ngspice\\\scriptsize + gm/ID LUTs};
    \node[agent3, left=of ng] (gds) {GDS\\\scriptsize gLayout / Magic};

    \draw[arrow]  (nl) -- (rec);
    \draw[arrow]  (rec) -- (adv);
    \draw[arrow]  (adv) -- (ar);
    \draw[arrow2] (ar) -- (ng);
    \draw[arrow2] (ng) -- (gds);
  \end{tikzpicture}
  \end{center}
  \vspace{0.3em}
  \begin{itemize}\footnotesize
    \item \cmd{analog-topology-recommender}: maps a one-line spec
          ("rail-to-rail OTA, 10 MHz GBW") to a registered topology
          (Miller OTA, telescopic, folded cascode, SAR-ADC,\ldots).
    \item \cmd{analog-sizing-advisor}: consumes the chosen topology plus
          gm/ID LUTs; proposes a starting W/L + bias point; hands off to
          \term{AutoresearchRunner} for the SPICE-scored greedy refinement.
    \item \cmd{gf180-docker-analog}: finalises layout + runs DRC/LVS
          inside \cmd{hpretl/iic-osic-tools}.
  \end{itemize}
  \note{%
    The analog loop is where \cmd{eda-agents} has the most depth. The
    methodology is \term{gm/ID}: pre-computed LUTs encode the
    device-level trade-off (transconductance vs current) for each
    transistor type on each PDK. The sizing advisor uses those tables to
    hand the autoresearch runner a good starting point -- not a random
    one -- so the subsequent SPICE budget is spent on refinement, not
    discovery.\par
    Supported circuit topologies today (registered as
    \cmd{CircuitTopology} subclasses under
    \cmd{src/eda\_agents/topologies/}): Miller OTA, AnalogAcademy OTA,
    GF180 OTA, StrongArm comparator, SAR ADC (8b, both transistor-level
    and behavioral variants). Topology-agnostic core means ``add another
    circuit'' is \cmd{\~200} lines of code per topology.\par
    Layout: \cmd{gLayout} (Python-based parametric layout) in a sibling
    venv, Magic for DRC/LVS, KLayout for streamout. The ADK harness
    drives those as sub-agents \cmd{FlowRunner}, \cmd{DRCChecker},
    \cmd{LVSVerifier}.}
\end{frame}

\begin{frame}{Key skills and artefacts -- analog}
  \footnotesize
  \begin{tabular}{@{}p{0.40\textwidth}p{0.55\textwidth}@{}}
    \toprule
    \textbf{Asset} & \textbf{What it does} \\
    \midrule
    \cmd{analog.gmid\_sizing} skill & Prompt template + policy for gm/ID
        starting-point selection. \\
    \cmd{analog.explorer} skill & Proposes next-step sizing moves; reads
        \cmd{program.md} and \cmd{results.tsv}. \\
    \cmd{analog.corner\_validator} & Re-runs the best candidate across
        PVT corners before accepting. \\
    \cmd{analog.behavioral\_primitives} & Catalogue of XSPICE + Verilog-A
        primitives (ideal comp, clock gen, sampler) for top-down SAR,
        PLL, LDO blocks. \\
    \cmd{analog.sar\_adc\_design} & End-to-end SAR ADC skill (8b-11b;
        references \cmd{sar\_adc\_8bit\_behavioral} topology). \\
    \midrule
    \cmd{AutoresearchRunner}
        (\cmd{src/.../autoresearch\_runner.py}) & The greedy loop.
        Persists \cmd{program.md} (narrative) + \cmd{results.tsv} for
        resumable runs. \\
    \cmd{SpiceRunner}
        (\cmd{src/.../spice\_runner.py}) & Sync + async ngspice with
        measurement parsing. \\
    \bottomrule
  \end{tabular}
  \note{%
    All the skill names are dotted namespace entries registered in
    \cmd{src/eda\_agents/skills/analog.py}. The MCP server exports them
    so any agent (ours or a user's) can request the prompt by name
    rather than copy-pasting text.\par
    The \cmd{program.md} + \cmd{results.tsv} pair is a Karpathy-style
    autoresearch trace: the LLM reads its own prior thinking and
    measurements, making runs \emph{resumable}. That is how we handle
    10+ hour budgets without blowing context windows.}
\end{frame}
